\section*{अध्याय \ref{ch:g11:trigo} --- त्रिकोणमिति: इकाई वृत्त}

\begin{solution}{exo:g11:trigo:1}
\omterm{def:g11:trigo:radian}{रेडियन} में: $30^\circ = \frac{\pi}{6}$, $45^\circ = \frac{\pi}{4}$, $120^\circ = \frac{2\pi}{3}$, $270^\circ = \frac{3\pi}{2}$।

अंशों में: $\frac{\pi}{5} = 36^\circ$, $\frac{5\pi}{6} = 150^\circ$, $\frac{7\pi}{4} = 315^\circ$, $\frac{2\pi}{9} = 40^\circ$।
\end{solution}

\begin{solution}{exo:g11:trigo:2}
$\frac{2\pi}{3}$ और $-\frac{\pi}{4}$ पहले से ही $\intoc{-\pi}{\pi}$ में हैं: क्रमशः दूसरा और चौथा चतुर्थांश।

$\frac{17\pi}{6} - 2\pi = \frac{5\pi}{6}$: दूसरा चतुर्थांश, ऋणात्मक $x$-अक्ष के पास।

$-\frac{7\pi}{2} + 4\pi = \frac{\pi}{2}$: वृत्त का शीर्ष, $(0, 1)$।
\end{solution}

\begin{solution}{exo:g11:trigo:3}
$\cos\frac{3\pi}{4} = \cos\left(\pi - \frac\pi4\right)
= -\cos\frac\pi4 = -\frac{\sqrt2}{2}$।

$\sin\frac{5\pi}{6} = \sin\left(\pi - \frac\pi6\right)
= \sin\frac\pi6 = \frac12$।

$\cos\left(-\frac\pi3\right) = \cos\frac\pi3 = \frac12$।

$\sin\frac{4\pi}{3} = \sin\left(\pi + \frac\pi3\right)
= -\sin\frac\pi3 = -\frac{\sqrt3}{2}$।

$\cos\frac{11\pi}{6} = \cos\left(-\frac\pi6 + 2\pi\right)
= \cos\frac\pi6 = \frac{\sqrt3}{2}$।
\end{solution}

\begin{solution}{exo:g11:trigo:4}
$\sin^2 t = 1 - \cos^2 t = 1 - \frac{25}{169} = \frac{144}{169}$, इसलिए $\sin t = \pm\frac{12}{13}$। $\intoo{-\frac\pi2}{0}$ पर (चौथा चतुर्थांश) कोटि ऋणात्मक है: $\sin t = -\frac{12}{13}$।
\end{solution}

\begin{solution}{exo:g11:trigo:5}
\cref{prop:g11:trigo:associated} का उपयोग करने पर:
\[
A = (-\cos x) + \cos x + \cos x - \cos x = 0 ,
\]
\[
B = \sin x + (-\sin x) + (-\sin x) = -\sin x .
\]
\end{solution}

\begin{solution}{exo:g11:trigo:6}
$\cos t = \frac{\sqrt3}{2}$: एक हल $\frac\pi6$ है, इसलिए $t = \pm\frac\pi6$ ($2k\pi$ के खिसकाव \omterm{def:g10:numbers:interval}{अंतराल} से बाहर ले जाते हैं)।

$\sin t = \frac12$: $t = \frac\pi6$ या $t = \pi - \frac\pi6 =
\frac{5\pi}{6}$।

$\cos t = -1$: \omterm{def:g10:coordgeom:system}{भुज} $-1$ वाला अकेला बिंदु $M(\pi)$ है, इसलिए $t = \pi$।
\end{solution}

\begin{solution}{exo:g11:trigo:7}
$\sin t = -\frac{\sqrt3}{2}$: $a = -\frac\pi3$ से परिवार $-\frac\pi3 + 2k\pi$ और $\pi + \frac\pi3 + 2k\pi$ हैं; $\intco{0}{2\pi}$ में: $t = \frac{4\pi}{3}$ और $t = \frac{5\pi}{3}$।

$\cos t = 0$: $t = \frac\pi2$ और $t = \frac{3\pi}{2}$।

$2\sin t + 1 = 0$ का अर्थ है $\sin t = -\frac12$: $a = -\frac\pi6$ से, $\intco{0}{2\pi}$ में: $t = \frac{7\pi}{6}$ और $t = \frac{11\pi}{6}$।
\end{solution}

\begin{solution}{exo:g11:trigo:8}
\cref{thm:g11:trigo:equations} से $\cos 2t = \cos t$ का अर्थ है $2t = t + 2k\pi$ या $2t = -t + 2k\pi$, अर्थात् $t = 2k\pi$ या $t = \frac{2k\pi}{3}$ ($k \in \Z$)। पहला
परिवार दूसरे में समा जाता है ($k$ को $3$ का गुणज लेने पर), इसलिए
हल-समुच्चय
\[
S = \left\{\frac{2k\pi}{3} : k \in \Z\right\}
\]
है --- वृत्त पर तीनों बिंदु $M(0)$, $M\!\left(\frac{2\pi}{3}\right)$, $M\!\left(\frac{4\pi}{3}\right)$।
\end{solution}

\begin{solution}{exo:g11:trigo:9}
दोनों वर्गों का प्रसार करके $\cos^2 x + \sin^2 x = 1$ का उपयोग करने पर:
\[
(\cos^2 x + 2\sin x\cos x + \sin^2 x)
+ (\cos^2 x - 2\sin x\cos x + \sin^2 x)
= 1 + 1 = 2 .
\]
\end{solution}

\begin{solution}{exo:g11:trigo:10}
बिना फिसले लुढ़कने का अर्थ है कि चाप-लंबाई तय की गई दूरी के बराबर है: $r\theta = 150$
cm से $\theta = \frac{150}{30} = 5$ rad $= 5 \times \frac{180}{\pi} \approx
286.5^\circ$ मिलता है। एक पूरे चक्कर में साइकिल परिधि $2\pi \times 30 = 60\pi \approx 188.5$ cm, यानी
लगभग $1.88$ m आगे बढ़ती है।
\end{solution}

\begin{solution}{exo:g11:trigo:11}
इकाई वृत्त पर ठीक $\frac12$ \omterm{def:g10:coordgeom:system}{भुज} वाले बिंदु $M\!\left(\pm\frac\pi3\right)$ हैं। \omterm{def:g10:coordgeom:system}{भुज} \emph{अधिक से अधिक} $\frac12$
वाले बिंदु ऊर्ध्वाधर रेखा $x = \frac12$ के बाईं ओर का चाप बनाते हैं, जिस पर $t$
$\frac\pi3$ से $\pi$ तक, या $-\pi$ से $-\frac\pi3$ तक जाने पर चला जाता है। अतः $\intoc{-\pi}{\pi}$ पर:
\[
S = \intoc{-\pi}{-\frac{\pi}{3}} \cup \intcc{\frac{\pi}{3}}{\pi}.
\]
(सिरे $\pm\frac\pi3$ सम्मिलित हैं, क्योंकि असमिका चौड़ी है।)
\end{solution}

\begin{solution}{pb:g11:trigo:1}
\textbf{1.} $30^\circ = \frac{\pi}{6}$; $135^\circ =
\frac{3\pi}{4}$; $300^\circ = \frac{5\pi}{3}$। और $\frac{3\pi}{4} = 135^\circ$; $\frac{5\pi}{6} = 150^\circ$।

\textbf{2.} लंबाई $r\,t = 3 \times 2 = 6$ m। \omterm{def:g11:trigo:radian}{रेडियन} चाप-लंबाई को एक सीधा गुणनफल बना देते हैं
--- अंश हर सूत्र में एक $\frac{\pi}{180}$ घसीट लाते।

\textbf{3.} घूमे हुए \omterm{def:g11:trigo:radian}{रेडियन} $= \frac{\text{दूरी}}
{\text{त्रिज्या}} = \frac{1000}{0.30} \approx 3\,333$ rad; $2\pi$ से भाग देने पर: लगभग $530$
पूरे चक्कर।

\textbf{4.} $\cos\frac{2\pi}{3} = -\frac12$; $\sin\left(-\frac\pi4\right) = -\frac{\sqrt2}{2}$; $\frac{19\pi}{6} - 2\pi = \frac{7\pi}{6}$, इसलिए $\cos\frac{19\pi}{6} = \cos\frac{7\pi}{6} =
-\frac{\sqrt3}{2}$।

\textbf{5.} $\cos^2 t = 1 - \frac{9}{25} = \frac{16}{25}$; दूसरे चतुर्थांश में \omterm{def:g11:trigo:cossin}{कोज्या} ऋणात्मक है: $\cos t = -\frac45$, और $\tan t = \frac{3/5}{-4/5} = -\frac34$।

\textbf{6.} $t$ मिनट में चक्र $\frac{t}{30} \times 2\pi = \frac{\pi t}{15}$ घूमता है। $t = 0$ पर: $h = 65 - 60 = 5$ m --- यानी
सबसे नीचे का सवारी-मंच (केंद्र की ऊँचाई घटा त्रिज्या)। ऋण चिह्न आपको तब
केंद्र से \emph{नीचे} रखता है जब घूमा हुआ कोण $0$ हो।

\textbf{7.} $h(7.5) = 65 - 60\cos\frac{\pi}{2} = 65$ m (केंद्र की ऊँचाई); $h(15) = 65 + 60 = 125$ m (शिखर); $h(20) = 65 - 60\cos\frac{4\pi}{3} = 65 + 30 = 95$ m।

\textbf{8.} $65 - 60\cos\frac{\pi t}{15} = 95$ से $\cos\frac{\pi t}{15} = -\frac12$ मिलता है: पहला हल $\frac{\pi t}{15} = \frac{2\pi}{3}$, अर्थात् $t = 10$ मिनट।

\textbf{9.} $\frac{\pi t}{15} \in \intcc{\frac{2\pi}{3}}{\frac{4\pi}{3}}$ के लिए $\cos\frac{\pi t}{15} \leq -\frac12$, अर्थात् $t \in \intcc{10}{20}$: सवारी के दस मिनट $95$ m से
ऊपर, जो $t = 15$ पर शिखर के परितः सममित हैं।

\textbf{10.} पहला मिनट:
\[
h(1) - h(0) = 60\left(1 - \cos\frac{\pi}{15}\right)
\approx 1.3 \text{ m}.
\]
मिनट $7$ और $8$ के बीच:
\[
60\left(\cos\frac{7\pi}{15} - \cos\frac{8\pi}{15}\right)
\approx 12.5 \text{ m}
\]
--- यानी लगभग दस गुना अधिक। ऊँचाई सबसे तेज़ी से तब बदलती है जब केंद्र के
स्तर से गुज़रा जाता है, और सबसे नीचे तथा सबसे ऊपर के पास वह लगभग ठहर जाती
है: ज्यावक्र अपने बीच में तीखा और अपने सिरों पर सपाट होता है।

\textbf{11.} त्रिज्या $= \frac{40 - 4}{2} = 18$ m, केंद्र $4 + 18 = 22$ m पर, आवर्तकाल $12$ मिनट:
$h(t) = 22 - 18\cos\left(\frac{\pi t}{6}\right)$।

\textbf{12.} $\sin\frac{\pi t}{6} = 1$ होने पर \omterm{def:g10:functions:extrema}{अधिकतम} $7 + 3 = 10$ m: पहली बार $t = 3$ पर (सुबह 3 बजे);
\omterm{def:g10:functions:extrema}{न्यूनतम} $4$ m पहली बार $t = 9$ पर; आवर्तकाल $\frac{2\pi}{\pi/6} = 12$ घंटे।

\textbf{13.} $\frac{\pi t}{6} \in \intcc{\frac\pi6}{\frac{5\pi}{6}}$ के लिए $\sin\frac{\pi t}{6} \geq \frac12$: $t \in \intcc{1}{5}$ --- यानी रात 1 बजे से सुबह 5 बजे तक
की चार घंटे की खिड़की, जो एक आवर्तकाल बाद $t \in \intcc{13}{17}$ पर फिर खुलती है (दोपहर 1
बजे से शाम 5 बजे तक)।

\textbf{14.} ऊर्ध्वाधर अक्ष का दर्पण कोण $t$ वाले बिंदु को कोण $\pi - t$ वाले
बिंदु पर भेजता है और कोटि बनाए रखता है: $\sin(\pi - t) = \sin t$। केंद्र के परितः आधा चक्कर
$t$ को $\pi + t$ पर भेजता है और दोनों \omterm{def:g10:coordgeom:system}{निर्देशांकों} के चिह्न बदल देता है:
$\cos(\pi + t) = -\cos t$। और विकर्ण का दर्पण \omterm{def:g10:coordgeom:system}{भुज} तथा कोटि आपस में बदल देता है: $\cos\left(\frac\pi2 - t\right) = \sin t$ --- \omterm{def:g10:proba:operations}{पूरक}
कोणों वाला वही ``को''।

\textbf{15.} $\sin(-t) = -\sin t$: विषम (उसके \omterm{def:g10:functions:graph}{आलेख} की अर्ध-घूर्णन सममिति); $\cos(-t) = \cos t$: सम
(दर्पण)। और $\sin t + t^3$ दो \omterm{def:g11:func:parity}{विषम फलनों} का योग है: विषम।

\textbf{16.} \omterm{def:g10:algebra:expand}{गुणनखंडन}: $(2\sin t + 1)(\sin t - 1) = 0$: $\sin t = 1$ से $t = \frac\pi2$ मिलता है; $\sin t = -\frac12$ से $t = \frac{7\pi}{6}$ और
$t = \frac{11\pi}{6}$। वृत्त पर तीन हल।

\textbf{17.} $1$ rad $= \frac{180}{\pi} \approx 57.3^\circ$। इसलिए $\sin(1) \approx 0.841$ जबकि $\sin(1^\circ) \approx
0.0175$: लगभग $50$ का गुणक।
सीख: किसी भी त्रिकोणमितीय पठन पर भरोसा करने से पहले विधा की बत्ती देख लीजिए।

\textbf{18.} \omterm{def:g10:coordgeom:system}{भुज} और कोटि विकर्ण पर बराबर होते हैं: $t = \frac\pi4$ और $t = \frac{5\pi}{4}$।

\textbf{19.} $\sin(0.1) = 0.09983\ldots$ बनाम $0.1$: आपेक्षिक त्रुटि लगभग $0.17\,\%$। छोटे कोणों के लिए
जीवा चाप से चिपकी रहती है --- और कक्षा 12 की सीमा $\frac{\sin t}{t} \to 1$ ठीक इसी प्रेक्षण को
प्रमेय बना देती है।

\textbf{20.} \omterm{def:g11:trigo:radian}{रेडियन}: कोण $=$ इकाई वृत्त पर चाप, इसलिए कोण लंबाइयों की तरह
गणना में आ जाते हैं। \omterm{def:g11:trigo:cossin}{कोज्या} और \omterm{def:g11:trigo:cossin}{ज्या}: घड़ी की सुई के \omterm{def:g10:coordgeom:system}{निर्देशांक} --- ज़मीन पर
छाया, दीवार पर ऊँचाई। वृत्त की सममितियाँ: सर्वसमिकाओं की पूरी सूची, जो
दर्पणों और अर्ध-घूर्णनों से पढ़ ली जाती है। \omterm{def:g10:algebra:equation}{समीकरण} $\cos t = c$: घड़ी की
समय-सारणी, जिसमें वृत्त हर हल एक साथ दिखा देता है। वेश: पहिये ने $65 - 60\cos$ पहना,
ज्वार ने $7 + 3\sin$ --- वही घड़ी, वही गणित।
\end{solution}
